\documentclass[12pt]{article}
\usepackage[utf8]{inputenc}
\usepackage[a4paper,margin=1in]{geometry}
\usepackage{amsmath, amssymb}
\usepackage{graphicx}
\usepackage{hyperref}
\usepackage{longtable}
\usepackage{booktabs}
\usepackage{algorithm}
\usepackage{algpseudocode}

\title{Two--Layer Weekly Service Dispatch:\\Scheduling and Routing for NES}
\author{VB}
\date{}

\begin{document}

\maketitle

% ===================================================================
\section*{Problem Description}
% ===================================================================

New England SteamWorks (NES) dispatches multiple technicians with service vehicles to perform maintenance jobs across a wide geographic region every week.
The planning task is decomposed into two sequential optimisation layers:

\begin{enumerate}
    \item \textbf{Layer~1 --- Weekly Scheduling:} Assign candidate jobs to (technician, vehicle, day) triples so as to maximise total weighted job value while keeping workloads approximately uniform across workers and vehicles.
    \item \textbf{Layer~2 --- Daily Routing:} For each (technician, vehicle, day) bundle produced by Layer~1, solve an intra--cluster Vehicle Routing Problem to determine the optimal visit sequence.
\end{enumerate}

The engine must:
\begin{itemize}
    \item Maximise the total weighted value of jobs served each week (Layer~1);
    \item Minimise total driving distance within each daily route (Layer~2);
    \item Respect capacity, time, skill, vehicle, and geographic constraints;
    \item Produce reviewable route/day proposals (not direct calendar writes).
\end{itemize}

% ===================================================================
\section*{Key Objectives (KPIs)}
% ===================================================================
\begin{itemize}
    \item \textbf{Primary KPI:} Weighted job score (combines revenue, aging pressure, and seasonal fairness).
    \item \textbf{Secondary KPIs:} Workload uniformity, total driving distance, route compactness, area readiness.
\end{itemize}

% ===================================================================
\section*{System Overview}
% ===================================================================

\begin{itemize}
    \item \textbf{ServiceM8} --- system of record for jobs, technicians, and operational data.
    \item \textbf{Airtable} --- scheduling rules, policy settings, review layers.
    \item \textbf{Python Engine} --- scheduling (Layer~1), clustering, routing (Layer~2), and output generation.
    \item \textbf{Make\,/\,Zapier} --- orchestration and automation between the above.
\end{itemize}

The Python engine consumes structured inputs and produces structured outputs; it is independently testable.

% ===================================================================
\section*{Sets and Indices}
% ===================================================================
\begin{itemize}
    \item $\mathcal{J} = \{1,\dots,J\}$ --- weekly candidate job pool (after exclusion filters).
    \item $\mathcal{J}^H = \{j \in \mathcal{J} : h_j = 1\}$ --- subset of two-man (helper-needed) jobs.
    \item $\mathcal{T} = \{1,\dots,T\}$ --- available technicians.
    \item $\mathcal{K} = \{1,\dots,K\}$ --- available vehicles.
    \item $\mathcal{D} = \{1,\dots,D\}$ --- working days in the planning week ($D \leq 5$).
    \item $\mathcal{A}$ --- set of geographic areas / zones.
    \item For a given technician $t \in \mathcal{T}$ on day $d \in \mathcal{D}$, the \emph{cluster} of feasible jobs is $\mathcal{C}_{t,d} \subseteq \mathcal{J}$ (Section~\ref{sec:clustering}).
\end{itemize}

% ===================================================================
\section*{Constraints and Parameters}
% ===================================================================

\subsection*{Vehicle Types}
\begin{itemize}
    \item \textbf{V-2} --- required for Radiator Runs (special route type).
    \item \textbf{V-4} --- weekly capacity switch (available / unavailable toggles capacity).
    \item \textbf{V-5, V-6, V-7} --- standard service vehicles.
\end{itemize}

Each vehicle $k$ has a set of capability tags $\mathit{cap}_k$, a daily availability vector, a speed $v_k$ (m/min), a cost rate $c_k$ (\$/m), and a stop capacity $Q_k$.

\subsection*{Technician Parameters}
Each technician $t$ is characterised by:
\begin{itemize}
    \item Skill set $\mathit{skill}_t$.
    \item Weekly availability after exceptions: $\mathcal{D}_t \subseteq \mathcal{D}$.
    \item Home depot / start location $(y_t, x_t)$.
\end{itemize}

\subsection*{Job Parameters}
Each job $j \in \mathcal{J}$ carries:
\begin{itemize}
    \item Address with geocoordinates $(\mathrm{lat}_j, \mathrm{lon}_j)$.
    \item Area label $a_j \in \mathcal{A}$.
    \item Required route type $\rho_j$ (normal, radiator, NH overnight, helper-needed).
    \item Helper flag $h_j \in \{0,1\}$.
    \item Age in days $\alpha_j$ (time since job was created).
    \item Base revenue / value $w_j$.
    \item Service time $\tau_j$ (minutes) --- average on-site duration.
\end{itemize}

\subsection*{Time and Distance}
\begin{itemize}
    \item Cluster radius: $R$ (meters) --- maximum Haversine distance from technician start to any job in the cluster.
    \item Maximum inter-stop distance: $r$ (meters) --- maximum distance between consecutive stops.
    \item Vehicle speed: $v_k$ (m/min) --- depends on vehicle $k \in \mathcal{K}$.
    \item Cost per metre driven: $c_k$ (\$/m) --- depends on vehicle $k \in \mathcal{K}$.
    \item Service time per stop: $\tau_j$ (minutes) --- depends on job $j \in \mathcal{J}$.
    \item Vehicle capacity: $Q_k$ --- maximum number of stops vehicle $k$ can serve per route.
    \item Maximum daily route duration: $T_{\max}$ (minutes).
    \item Maximum stops per route: $P$ (global upper bound; vehicle capacity $Q_k \leq P$).
\end{itemize}

\subsection*{Seasonal Weights}
Scoring weights vary by season:
\[
\text{score}_j \;=\; w_g \cdot \text{geo}_j \;+\; w_f \cdot \text{fair}_j \;+\; w_a \cdot \text{age}_j \;+\; w_r \cdot \text{readiness}_j
\]
where $(w_g, w_f, w_a, w_r)$ shift between a geography--dominant profile (Mar--Sep) and a fairness--dominant profile (Oct--Feb).

\subsection*{Weekly Exceptions}
Before route construction the engine applies weekly exceptions:
\begin{itemize}
    \item Technician unavailability days.
    \item Vehicle unavailability.
    \item Pre-scheduled long jobs and boiler installs that consume partial capacity.
\end{itemize}
These reduce effective weekly capacity before the scheduling layer begins.

% ===================================================================
\section*{Input Data}
% ===================================================================

\subsection*{Job Data (\texttt{example\_jobs.csv})}
A structured export from ServiceM8 providing, for each candidate job:

{\small
\begin{longtable}{ccccccccccc}
\toprule
\textbf{ID} & \textbf{Address} & \textbf{Lat} & \textbf{Lon} & \textbf{Area} & \textbf{Route Type} & \textbf{Helper} & \textbf{Age} & \textbf{Value} & \textbf{Svc (min)} & \textbf{Status} \\
\midrule
J-001 & 123 Main St, Concord NH & 43.2081 & $-$71.5376 & Area-1 & normal & 0 & 14 & 250 & 30 & candidate \\
J-002 & 456 Elm St, Manchester NH & 42.9956 & $-$71.4548 & Area-2 & normal & 0 & 7  & 180 & 25 & candidate \\
J-003 & 789 Broadway, Nashua NH & 42.7654 & $-$71.4676 & Area-3 & normal & 0 & 21 & 320 & 45 & candidate \\
\bottomrule
\end{longtable}
}

\noindent Columns:
\begin{itemize}
    \item \texttt{job\_id} --- unique identifier.
    \item \texttt{address} --- street address (geocoded to lat/lon).
    \item \texttt{latitude}, \texttt{longitude} --- WGS-84 coordinates.
    \item \texttt{area\_id} --- geographic area label ($a_j \in \mathcal{A}$).
    \item \texttt{route\_type} --- \texttt{normal}, \texttt{radiator}, \texttt{nh\_overnight}, or \texttt{helper}.
    \item \texttt{helper\_needed} --- binary flag ($h_j$); 1 if a second person is required.
    \item \texttt{age\_days} --- days since the job was created ($\alpha_j$).
    \item \texttt{value} --- base revenue / job weight ($w_j$, \$).
    \item \texttt{service\_time\_min} --- expected on-site duration ($\tau_j$, minutes).
    \item \texttt{status} --- \texttt{candidate} (eligible), \texttt{excluded}, or \texttt{assigned}.
    \item \texttt{description} --- free-text note (not used by the optimiser).
\end{itemize}

The engine pre-computes a distance matrix $d_{i,j}$ for every pair of jobs using map-based shortest-path distances (OSMnx / NetworkX), falling back to Haversine when no road path exists.

\subsection*{Technician Data (\texttt{example\_technicians.csv})}
Technician roster with skill sets and home depot coordinates:

{\small
\begin{longtable}{cccccc}
\toprule
\textbf{ID} & \textbf{Name} & \textbf{Skills} & \textbf{Home Lat} & \textbf{Home Lon} & \textbf{Available Days} \\
\midrule
T-01 & Mike Sullivan  & boiler, radiator, \dots  & 43.2081 & $-$71.5376 & Mon--Fri \\
T-02 & Dave Patel     & boiler, steam\_trap, \dots & 42.9956 & $-$71.4548 & Mon--Fri \\
T-03 & Chris O'Brien  & boiler, radiator, \dots  & 43.0718 & $-$70.7626 & Mon--Thu \\
T-04 & Sam Torres     & boiler, steam\_trap, \dots & 42.7654 & $-$71.4676 & Mon--Fri \\
\bottomrule
\end{longtable}
}

\noindent Columns:
\begin{itemize}
    \item \texttt{tech\_id} --- unique technician identifier.
    \item \texttt{name} --- display name.
    \item \texttt{skills} --- comma-separated list of skill tags ($\mathit{skill}_t$).
    \item \texttt{home\_lat}, \texttt{home\_lon} --- home depot / start location ($y_t, x_t$).
    \item \texttt{available\_days} --- base weekly availability ($\mathcal{D}_t$ before exceptions).
\end{itemize}

\subsection*{Vehicle Data (\texttt{example\_vehicles.csv})}
Vehicle fleet with capability tags, speed, cost, and capacity:

{\small
\begin{longtable}{ccccccc}
\toprule
\textbf{ID} & \textbf{Type} & \textbf{Capabilities} & \textbf{Available Days} & \textbf{Speed (m/min)} & \textbf{Cost (\$/m)} & \textbf{Capacity} \\
\midrule
V-2 & radiator\_van     & radiator, install, \dots & Mon--Fri & 400 & 0.00045 & 8  \\
V-4 & capacity\_switch   & boiler, steam\_trap      & Mon, Wed, Fri & 500 & 0.00030 & 10 \\
V-5 & service\_truck     & boiler, steam\_trap, \dots & Mon--Fri & 500 & 0.00035 & 14 \\
V-6 & service\_truck     & boiler, steam\_trap, \dots & Mon--Fri & 500 & 0.00035 & 14 \\
V-7 & service\_truck     & boiler, radiator, \dots  & Mon--Fri & 450 & 0.00040 & 12 \\
\bottomrule
\end{longtable}
}

\noindent Columns:
\begin{itemize}
    \item \texttt{vehicle\_id} --- unique vehicle identifier.
    \item \texttt{vehicle\_type} --- vehicle class name.
    \item \texttt{capability\_tags} --- comma-separated list of route-type capabilities ($\mathit{cap}_k$).
    \item \texttt{available\_days} --- base weekly availability ($\mathcal{D}_k$ before exceptions).
    \item \texttt{speed\_mpm} --- driving speed $v_k$ (metres per minute).
    \item \texttt{cost\_per\_metre} --- cost rate $c_k$ (\$/metre).
    \item \texttt{capacity} --- maximum stops per route $Q_k$.
\end{itemize}

\subsection*{Weekly Exceptions Data (\texttt{example\_weekly\_exceptions.csv})}
Per-week overrides that restrict technician or vehicle availability:

{\small
\begin{longtable}{cccccc}
\toprule
\textbf{Exception ID} & \textbf{Scope Type} & \textbf{Scope ID} & \textbf{Day} & \textbf{Effect Type} & \textbf{Effect Value} \\
\midrule
EX-01 & technician & T-03 & Fri & unavailable & PTO \\
EX-02 & vehicle    & V-4  & Tue & unavailable & maintenance \\
EX-03 & vehicle    & V-4  & Thu & unavailable & maintenance \\
EX-04 & technician & T-01 & Wed & partial     & boiler\_install\_3hrs \\
EX-05 & technician & T-02 & Mon & unavailable & training \\
\bottomrule
\end{longtable}
}

\noindent Columns:
\begin{itemize}
    \item \texttt{exception\_id} --- unique identifier for the exception.
    \item \texttt{scope\_type} --- \texttt{technician} or \texttt{vehicle}.
    \item \texttt{scope\_id} --- the affected technician or vehicle ID.
    \item \texttt{day} --- the day of the week the exception applies to.
    \item \texttt{effect\_type} --- \texttt{unavailable} (entire day blocked) or \texttt{partial} (reduced capacity).
    \item \texttt{effect\_value} --- reason or detail (e.g.\ PTO, maintenance, pre-booked job).
\end{itemize}

Exceptions are applied during data loading to compute effective availability sets $\mathcal{D}_t$ and $\mathcal{D}_k$ and are consumed by the Layer~1 eligibility pre-filter.

\subsection*{Duplicate--Address Handling}
Jobs sharing the same address are grouped for route-shape evaluation (to avoid inflating stop density) and flagged for review. They are never silently removed.

% ===================================================================
\section*{Planning Pipeline (Two--Layer Architecture)}
\label{sec:pipeline}
% ===================================================================

The weekly planning pipeline runs the following stages in order:

\begin{enumerate}
    \item \textbf{Load \& Normalise} --- Import candidate pool, technician roster, vehicle roster, exceptions.
    \item \textbf{Exclude} --- Remove jobs already scheduled, cancelled, urgent/manual, or in-workflow.
    \item \textbf{Apply Weekly Exceptions} --- Compute true weekly capacity per technician--vehicle pair.
    \item \textbf{Reserve Special Routes} --- Place Radiator Runs (require V-2), NH Overnight Runs, and helper-needed routes \emph{before} normal filling.
    \item \textbf{Layer~1: Schedule (Section~\ref{sec:scheduling})} --- Assign remaining jobs to (technician, vehicle, day) triples via the scheduling MILP.
    \item \textbf{Layer~2: Cluster \& Route (Sections~\ref{sec:clustering}--\ref{sec:model})} --- For each (technician, vehicle, day) bundle, build a geographic cluster, then solve the intra--cluster VRP.
    \item \textbf{Select Standby Jobs} --- Choose plausible backup jobs per route.
    \item \textbf{Compute Area Readiness} --- Derive area communication layer (\texttt{Good} / \texttt{Moderate} / \texttt{Lean}).
    \item \textbf{Render Outputs} --- Route maps, flags, summaries, reason codes.
\end{enumerate}

% ===================================================================
\section*{Layer~1: Weekly Scheduling Model (MILP)}
\label{sec:scheduling}
% ===================================================================

The scheduling layer decides \emph{which} jobs to assign and \emph{to whom} (technician $t$, vehicle $k$, day $d$), seeking to maximise total weighted job value subject to eligibility, capacity, and workload-balance constraints.

\subsection*{Eligibility Pre-filter}

Define the binary parameter:
\[
e_{j,t,k,d} =
\begin{cases}
1 & \text{if job $j$ may be assigned to technician $t$ with vehicle $k$ on day $d$},\\
0 & \text{otherwise}.
\end{cases}
\]
$e_{j,t,k,d} = 1$ requires all of the following:
\begin{itemize}
    \item Technician $t$ is available on day $d$: $d \in \mathcal{D}_t$.
    \item Vehicle $k$ is available on day $d$: $d \in \mathcal{D}_k$.
    \item $\mathit{skill}_t \supseteq \text{req}(j)$: technician skills cover job requirements.
    \item $\mathit{cap}_k \supseteq \rho_j$: vehicle capabilities cover the route type.
    \item $\text{haversine}(y_t, x_t, \mathrm{lat}_j, \mathrm{lon}_j) \leq R$: job within cluster radius.
    \item No prohibited technician--job pairing.
\end{itemize}

\subsection*{Helper Eligibility Pre-filter}

For two-man jobs $j \in \mathcal{J}^H$ we define a helper-eligibility parameter:
\[
e^H_{j,t,d} =
\begin{cases}
1 & \text{if technician $t$ may assist job $j$ on day $d$},\\
0 & \text{otherwise}.
\end{cases}
\]
$e^H_{j,t,d} = 1$ requires:
\begin{itemize}
    \item $d \in \mathcal{D}_t$ (technician available on that day).
    \item $\mathit{skill}_t \supseteq \text{req}(j)$ (skills cover the job).
    \item $\text{haversine}(y_t, x_t, \mathrm{lat}_j, \mathrm{lon}_j) \leq R$ (job within reach).
\end{itemize}
Note that helper eligibility is independent of vehicle assignment; the helper travels with the primary technician's vehicle.

\subsection*{Decision Variables}
\[
y_{j,t,k,d} \in \{0,1\} \quad \forall\, j \in \mathcal{J},\; t \in \mathcal{T},\; k \in \mathcal{K},\; d \in \mathcal{D}
\]
$y_{j,t,k,d} = 1$ means job $j$ is assigned to technician $t$, vehicle $k$, on day $d$.

\[
y^H_{j,t,d} \in \{0,1\} \quad \forall\, j \in \mathcal{J}^H,\; t \in \mathcal{T},\; d \in \mathcal{D}
\]
$y^H_{j,t,d} = 1$ means technician $t$ serves as the \emph{helper} (second person) for two-man job $j$ on day $d$.

\subsection*{Auxiliary Variables}
\begin{itemize}
    \item $L_{t,d}$: total service-time load of technician $t$ on day $d$ (continuous, $\geq 0$).
    \item $L_{k,d}$: total service-time load of vehicle $k$ on day $d$ (continuous, $\geq 0$).
    \item $\bar{L}$: average per-slot load across all active (technician, day) pairs (parameter computed in pre-processing or auxiliary variable).
    \item $s^+_{t,d},\; s^-_{t,d} \geq 0$: positive/negative deviation of $L_{t,d}$ from $\bar{L}$ (for linearised uniformity).
\end{itemize}

\subsection*{Objective Function}

Maximise total weighted job value minus a uniformity penalty:
\begin{equation}
\max \;\;
\underbrace{\sum_{j \in \mathcal{J}} \sum_{t \in \mathcal{T}} \sum_{k \in \mathcal{K}} \sum_{d \in \mathcal{D}} w_j \, y_{j,t,k,d}}_{\text{total assigned weight}}
\;-\;
\underbrace{\lambda \sum_{t \in \mathcal{T}} \sum_{d \in \mathcal{D}} \bigl(s^+_{t,d} + s^-_{t,d}\bigr)}_{\text{uniformity penalty}}
\label{eq:sched_obj}
\end{equation}

where $\lambda \geq 0$ is a tunable penalty weight (larger $\lambda$ enforces stronger balance).

\subsection*{Subject to:}

\begin{align}
%--- (S1) each job assigned at most once ---
& \sum_{t \in \mathcal{T}} \sum_{k \in \mathcal{K}} \sum_{d \in \mathcal{D}} y_{j,t,k,d} \leq 1, \quad \forall\, j \in \mathcal{J}
    && \text{(Each job assigned at most once)} \label{eq:assign_once}\\[4pt]
%--- (S2) eligibility ---
& y_{j,t,k,d} \leq e_{j,t,k,d}, \quad \forall\, j,t,k,d
    && \text{(Eligibility)} \label{eq:eligible}\\[4pt]
%--- (S3) vehicle capacity per day ---
& \sum_{j \in \mathcal{J}} y_{j,t,k,d} \leq Q_k, \quad \forall\, t \in \mathcal{T},\; k \in \mathcal{K},\; d \in \mathcal{D}
    && \text{(Vehicle capacity per day)} \label{eq:sched_cap}\\[4pt]
%--- (S4) daily time budget ---
& \sum_{j \in \mathcal{J}} \tau_j \, y_{j,t,k,d} \leq T_{\max}, \quad \forall\, t,k,d
    && \text{(Service-time budget; travel added in Layer~2)} \label{eq:sched_time}\\[4pt]
%--- (S5) one vehicle per technician per day ---
& \sum_{k \in \mathcal{K}} \sum_{j \in \mathcal{J}} y_{j,t,k,d} \leq P \cdot z_{t,k,d}, \quad \forall\, t,d
    && \text{(Link to vehicle indicator)} \label{eq:one_veh_link}\\[4pt]
& \sum_{k \in \mathcal{K}} z_{t,k,d} \leq 1, \quad \forall\, t \in \mathcal{T},\; d \in \mathcal{D}
    && \text{(At most one vehicle per tech per day)} \label{eq:one_veh}\\[4pt]
& z_{t,k,d} \in \{0,1\}
    && \label{eq:z_bin}\\[4pt]
%--- (S6) one technician per vehicle per day ---
& \sum_{t \in \mathcal{T}} z_{t,k,d} \leq 1, \quad \forall\, k \in \mathcal{K},\; d \in \mathcal{D}
    && \text{(At most one tech per vehicle per day)} \label{eq:one_tech}\\[4pt]
%--- (S7) technician load definition (includes helper time) ---
& L_{t,d} = \sum_{j \in \mathcal{J}} \sum_{k \in \mathcal{K}} \tau_j \, y_{j,t,k,d} \;+\; \sum_{j \in \mathcal{J}^H} \tau_j \, y^H_{j,t,d}, \quad \forall\, t,d
    && \text{(Technician-day load incl.\ helper)} \label{eq:tech_load}\\[4pt]
%--- (S8) uniformity soft constraint ---
& L_{t,d} - \bar{L} = s^+_{t,d} - s^-_{t,d}, \quad \forall\, t,d
    && \text{(Deviation from mean load)} \label{eq:uniformity}\\[4pt]
& s^+_{t,d},\; s^-_{t,d} \geq 0
    && \label{eq:s_nonneg}\\[6pt]
%--- (S9) helper required for two-man jobs ---
& \sum_{t \in \mathcal{T}} y^H_{j,t,d} = \sum_{t \in \mathcal{T}} \sum_{k \in \mathcal{K}} y_{j,t,k,d}, \quad \forall\, j \in \mathcal{J}^H,\; d \in \mathcal{D}
    && \text{(Each scheduled helper job gets one helper)} \label{eq:helper_assign}\\[4pt]
%--- (S10) helper differs from primary technician ---
& y^H_{j,t,d} + \sum_{k \in \mathcal{K}} y_{j,t,k,d} \leq 1, \quad \forall\, j \in \mathcal{J}^H,\; t \in \mathcal{T},\; d \in \mathcal{D}
    && \text{(Helper $\neq$ primary technician)} \label{eq:helper_diff}\\[4pt]
%--- (S11) helper eligibility ---
& y^H_{j,t,d} \leq e^H_{j,t,d}, \quad \forall\, j \in \mathcal{J}^H,\; t \in \mathcal{T},\; d \in \mathcal{D}
    && \text{(Helper eligibility)} \label{eq:helper_elig}
\end{align}

\paragraph{Notes.}
\begin{itemize}
    \item Constraint~\eqref{eq:sched_time} uses service time only (no travel); travel feasibility is enforced in Layer~2 and may cause a small fraction of scheduled jobs to be dropped during routing.
    \item $\bar{L}$ may be fixed as $\displaystyle \bar{L} = \frac{\sum_{j} \tau_j}{|\{(t,d) : d \in \mathcal{D}_t\}|}$ (total service divided by active slots) or treated as a free variable alongside the deviations.
    \item For vehicle-level uniformity add analogous $s^+_{k,d}, s^-_{k,d}$ terms.
\end{itemize}

% ===================================================================
\section*{Layer~2: Clustering Phase}
\label{sec:clustering}
% ===================================================================

Given the Layer~1 assignment $(t, k, d) \mapsto \mathcal{J}_{t,k,d}$, the routing layer operates on each such bundle independently.  We adopt a \textbf{Cluster--first, Route--second} decomposition following the approach implemented in SW implementation.

\begin{algorithm}[H]
\caption{Geographic Cluster Construction for Technician $t$, Day $d$}
\begin{algorithmic}[1]
\Require Start location $(y_t, x_t)$, candidate jobs $\mathcal{J}$, cluster radius $R$, min cluster size $N_{\min}$, radius increment $\Delta R$, max iterations $I_{\max}$
\Ensure Cluster $\mathcal{C}_{t,d}$
\State $R_{\text{cur}} \gets R$
\For{$\text{iter} = 1$ \textbf{to} $I_{\max}$}
    \State $\mathcal{C}_{t,d} \gets \emptyset$
    \For{each $j \in \mathcal{J}$}
        \If{$\text{haversine}(y_t, x_t, \mathrm{lat}_j, \mathrm{lon}_j) \leq R_{\text{cur}}$
            \textbf{and} $\text{area}(j)$ compatible
            \textbf{and} $\text{skill}_t \supseteq \text{req}(j)$
            \textbf{and} vehicle compatible}
            \State $\mathcal{C}_{t,d} \gets \mathcal{C}_{t,d} \cup \{j\}$
        \EndIf
    \EndFor
    \If{$|\mathcal{C}_{t,d}| \geq N_{\min}$} \textbf{break}
    \Else \hspace{1em} $R_{\text{cur}} \gets R_{\text{cur}} + \Delta R$
    \EndIf
\EndFor
\State Build distance sub-matrix $d_{i,j}$ for $i,j \in \{0\} \cup \mathcal{C}_{t,d}$ using road network
\end{algorithmic}
\end{algorithm}

Area boundaries act as guidance, not hard walls: neighbouring areas may be mixed when geography meaningfully improves route quality.

% ===================================================================
\section*{Layer~2: Intra--Cluster Routing Model (MILP)}
\label{sec:model}
% ===================================================================

For each cluster $\mathcal{C}_{t,d}$ produced from the jobs assigned by Layer~1, we solve an Orienteering-style VRP as a Mixed-Integer Linear Program. Let $N = |\mathcal{C}_{t,d}|$ and index the depot (technician start) as node~$0$.

\subsection*{Parameters}
\begin{itemize}
\item $N$ nodes (job locations in the cluster) forming set $W = \{1,\dots,N\}$.
\item Depot node $0$ with distances $d_{0,i}$, $d_{i,0}$ for $i \in W$.
\item Road--network distances $d_{i,j}$ between jobs.
\item Cluster radius $R$ and inter-stop radius $r$ (meters).
\item Vehicle speed $v_k$ (m/min) and cost per metre $c_k$ (\$/m) for the assigned vehicle $k$.
\item Service time $\tau_j$ (minutes) for each job $j$.
\item Vehicle capacity $Q_k$ --- maximum stops for vehicle $k$ ($Q_k \leq P$).
\item Job score $w_j$ (revenue adjusted by seasonal weights).
\item Maximum route duration $T_{\max}$ (minutes).
\item Maximum number of stops $P$.
\end{itemize}

\subsection*{Decision Variables}
\begin{itemize}
    \item $x_{i,j} \in \{0,1\}$: 1 if the route traverses arc $(i,j)$, $0 \leq i,j \leq N$.
    \item $u_i \in \{0,1,\dots,P\}$: visit order (Miller--Tucker--Zemlin), $0 \leq i \leq N$.
\end{itemize}

\subsection*{Objective Function}

Maximize the total weighted score minus travel cost for the route $\Omega$ using vehicle $k$:
\[
\max \;\; \sum_{i=0}^{N}\sum_{\substack{j=0\\j\neq i}}^{N} w_{j}\, x_{i,j} \;-\; \sum_{i=0}^{N}\sum_{\substack{j=0\\j\neq i}}^{N} c_k \cdot d_{i,j} \cdot x_{i,j}
\]

\subsection*{Subject to:}

\begin{align}
& \sum_{i=1}^{N} \sum_{\substack{j=1\\j\neq i}}^{N} x_{i,j} \leq Q_k
    && \text{(Vehicle capacity)} \label{eq:capacity}\\[4pt]
& \sum_{i=0}^{N}\sum_{\substack{j=0\\j\neq i}}^{N} x_{i,j}\!\left(\frac{d_{i,j}}{v_k} + \tau_{j}\right) \leq T_{\max}
    && \text{(Daily time budget)} \label{eq:time}\\[4pt]
& d_{0,i} \leq R,\;\; d_{i,0} \leq R, \quad i \in W
    && \text{(Cluster radius from depot)} \label{eq:cluster}\\[4pt]
& x_{i,i} = 0, \quad i=0,\dots,N
    && \text{(No self-loops)} \label{eq:noloop}\\[4pt]
& \sum_{i=0}^{N} x_{i,j} \leq 1, \quad j=0,\dots,N
    && \text{(Enter each node at most once)} \label{eq:inonce}\\[4pt]
& \sum_{j=0}^{N} x_{i,j} \leq 1, \quad i=0,\dots,N
    && \text{(Leave each node at most once)} \label{eq:outonce}\\[4pt]
& \sum_{j=1}^{N} x_{0,j} = 1
    && \text{(Depart from depot)} \label{eq:depart}\\[4pt]
& \sum_{i=1}^{N} x_{i,0} = 1
    && \text{(Return to depot)} \label{eq:return}\\[4pt]
& u_i - u_j + P\, x_{i,j} \leq P - 1,\; 1 \leq i,j \leq N
    && \text{(MTZ subtour elimination)} \label{eq:mtz}\\[4pt]
& \sum_{j=0}^{N} x_{i,j} = \sum_{l=0}^{N} x_{l,i}, \quad i=0,\dots,N
    && \text{(Flow conservation)} \label{eq:flow}\\[4pt]
& d_{i,j} \leq r \;\;\text{or}\;\; x_{i,j}=0, \quad 1\leq i,j \leq N
    && \text{(Inter-stop distance limit)} \label{eq:interstop}
\end{align}

Constraints \eqref{eq:cluster} and \eqref{eq:interstop} are pre-processing filters: arcs violating them are removed from the feasible arc set before the solver runs, reducing the problem dimension (as done in the optimised MILP implementation).

% ===================================================================
\section*{Hard Constraints (Eligibility Layer)}
% ===================================================================

Before the MILP is invoked, jobs and assignments are filtered by hard constraints:

\begin{itemize}
    \item Job must be in the valid weekly candidate pool.
    \item Technician must be available on the planned day (after weekly exceptions).
    \item Vehicle must be available and capability-compatible with the route type.
    \item Special route prerequisites must be satisfied (e.g.\ V-2 for Radiator Runs).
    \item Helper-needed jobs require adequate staffing.
    \item Prohibited technician--job pairings are excluded.
\end{itemize}

Each exclusion is annotated with a reason code for traceability.

% ===================================================================
\section*{Special Route Types}
% ===================================================================

Special routes are planned \emph{before} normal route filling to reserve capacity:

\begin{enumerate}
    \item \textbf{Radiator Runs} --- Require vehicle V-2. Clustered separately.
    \item \textbf{NH Overnight Runs} --- Long-distance routes requiring overnight stay.
    \item \textbf{Two-Man / Helper Routes} --- Jobs flagged with $h_j = 1$; need two technicians or a helper.
\end{enumerate}

After special routes consume capacity, the remaining capacity feeds normal route construction.

% ===================================================================
\section*{Solution Methods}
% ===================================================================

For each cluster the engine can apply one or more of the following, matching the algorithms in the codebase:

\begin{enumerate}
    \item \textbf{MILP} (CBC solver via Pyomo) --- exact solution for small--to--medium clusters.
    \item \textbf{Simulated Annealing} --- metaheuristic for larger clusters; supports street-filling variant.
    \item \textbf{Greedy Heuristic} --- fast constructive heuristic with revenue-to-distance scoring.
    \item \textbf{Brute-Force Search} --- exhaustive enumeration for very small instances.
\end{enumerate}

A time budget of 120~seconds is enforced across all methods.

% ===================================================================
\section*{Standby Job Selection}
% ===================================================================

After route construction, unassigned but feasible jobs are ranked as standby replacements for each route. Standby jobs are output separately from assigned jobs so they can be deployed if a scheduled job drops out.

% ===================================================================
\section*{Area Readiness}
% ===================================================================

The engine computes an area--readiness indicator for each geographic area $a \in \mathcal{A}$:

\[
\text{Readiness}(a) \in \{\texttt{Good},\;\texttt{Moderate},\;\texttt{Lean}\}
\]

based on the density and quality of planned activity in that area. This is a communication layer for customer--facing scheduling messages and does not drive route construction.

% ===================================================================
\section*{Post-Route Update Rule}
% ===================================================================

After a weekly plan is finalised and approved, assigned jobs $j$ with $\sum_{i} x_{i,j} = 1$ are removed from the candidate pool for subsequent planning iterations.

% ===================================================================
\section*{Output Artefacts}
% ===================================================================

Each weekly run produces:
\begin{itemize}
    \item Proposed route/day assignments per technician.
    \item Standby job list per route.
    \item Plotted map for each route/day.
    \item Area readiness values and plain-text lookup block.
    \item Review flags: duplicate-address groups, vehicle bottlenecks, cross-area blends, weak standby pools.
    \item Reason codes traceable to the rules that produced each decision.
\end{itemize}

Outputs are designed for human review by a non-technical owner; the engine does \emph{not} auto-book jobs.

% ===================================================================
\section*{Owner--Controlled Policy Inputs}
% ===================================================================

The optimisation engine exposes a compact set of tunable parameters that a non-technical business owner can adjust---via Airtable fields or a configuration file---to steer planning behaviour without modifying code.

\paragraph{Workload-balance weight ($\lambda$).}
Controls the trade-off between maximising total job value and equalising daily loads across technicians (Eq.~\ref{eq:sched_obj}).
Setting $\lambda = 0$ yields a pure value-maximising schedule; increasing $\lambda$ smooths the workload at the cost of potentially leaving some lower-value jobs unassigned.

\paragraph{Maximum daily route duration ($T_{\max}$).}
Caps the total travel-plus-service time for any single route.
Lowering $T_{\max}$ produces shorter days and leaves more buffer for delays; raising it allows more jobs per day but increases overtime risk.

\paragraph{Maximum stops per route ($P$) and vehicle capacities ($Q_k$).}
$P$ is a global upper bound; each vehicle's capacity $Q_k \leq P$ can be set independently.
Reducing $Q_k$ for a specific vehicle reserves slack for ad-hoc stops or equipment space.

\paragraph{Cluster radius ($R$) and inter-stop limit ($r$).}
$R$ defines how far from a technician's home depot jobs may be included; $r$ limits the maximum gap between consecutive stops.
Tightening these values produces geographically compact routes with shorter drive legs; relaxing them allows the engine to reach more jobs at the expense of longer routes.

\paragraph{Seasonal weight profile $(w_g, w_f, w_a, w_r)$.}
Four weights that shift the job-scoring formula between a geography-dominant mode (Mar--Sep) and a fairness-dominant mode (Oct--Feb).
The owner sets two seasonal profiles; the engine selects the active profile based on the planning-week date.

\paragraph{Weekly exceptions.}
Technician and vehicle unavailability entries (PTO, maintenance, pre-booked installs) entered each week.
These directly shrink the effective capacity seen by Layer~1 and must be finalised before the planning run.

\paragraph{Special-route reservations.}
The owner marks which jobs require Radiator Run (V-2), NH Overnight, or helper-needed handling.
Reserved routes are locked before normal filling, so the owner controls how much capacity is set aside for these categories.

\paragraph{Prohibited pairings and area definitions.}
Technician--job exclusions (e.g.\ customer-requested) and the geographic area map ($\mathcal{A}$) are maintained in Airtable.
Changes take effect on the next planning run without engine modification.

% ===================================================================
\section*{Major Risks and Ambiguities}
% ===================================================================

\paragraph{Layer-1 / Layer-2 feasibility gap.}
The scheduling MILP (Layer~1) budgets service time only (Eq.~\ref{eq:sched_time}); it does not account for travel time between stops.
Consequently, some bundles passed to Layer~2 may exceed $T_{\max}$ once routing distances are added, forcing jobs to be dropped at the routing stage.
Until an iterative Layer-1\,/\,Layer-2 feedback loop is implemented, a conservative $T_{\max}$ buffer (e.g.\ 80\,\% of the true day length) should be used in Layer~1.

\paragraph{Distance-matrix quality.}
Road-network distances are computed via OSMnx/NetworkX, with a Haversine fallback when no path exists.
If the underlying OpenStreetMap data is incomplete or stale for rural New England roads, route durations may be under- or over-estimated.
Periodic comparison with Google Maps drive-time queries is recommended.

\paragraph{Solver scalability.}
The Layer~1 MILP has $|\mathcal{J}| \times |\mathcal{T}| \times |\mathcal{K}| \times |\mathcal{D}|$ binary variables.
For the current fleet (30 jobs, 4 technicians, 5 vehicles, 5 days = 3\,000 variables) the problem is tractable, but scaling to 200+ jobs or 10+ technicians may require solver time-limits, column generation, or a decomposition heuristic.

\paragraph{Geocoding accuracy.}
Job coordinates originate from ServiceM8 address geocoding.
Incorrect or ambiguous addresses produce wrong cluster assignments and distance estimates.
A pre-run validation step that flags coordinates outside known service territory is advisable.

\paragraph{Exception staleness.}
Weekly exceptions must be entered before the planning run.
Late-arriving changes (e.g.\ same-day sick leave) are not captured and require manual route adjustment or a re-run, which the current pipeline does not automate.

\paragraph{Seasonal-weight calibration.}
The scoring weights $(w_g, w_f, w_a, w_r)$ are currently set by judgement rather than historical data.
Until job-outcome data is available for back-testing, the weights should be reviewed periodically and adjusted if dispatchers consistently override engine recommendations.

\paragraph{Helper-job coordination.}
Two-man jobs ($h_j = 1$) are now modelled explicitly in Layer~1 via the helper variable $y^H_{j,t,d}$ and constraints \eqref{eq:helper_assign}--\eqref{eq:helper_elig}.
However, the model assumes the helper travels with the primary technician's vehicle and does not account for the helper's own travel time to the job site.
If helpers must drive separately, an additional travel-feasibility check is needed.

\paragraph{Single-week planning horizon.}
Each run optimises one week in isolation.
Jobs deferred this week re-enter the pool next week with a higher age score, but there is no explicit multi-week look-ahead to prevent systematic starvation of low-value areas.

% ===================================================================
\section*{Recommended Engine Structure}
\label{sec:engine_structure}
% ===================================================================

The Python engine is organised as a pipeline of self-contained stages, each with explicit inputs, outputs, and a single responsibility.
The recommended package layout is:

\begin{verbatim}
nes_dispatch/
    __init__.py
    data/
        loaders.py          # CSV / Airtable readers -> dataclasses
        models.py           # Job, Technician, Vehicle, WeeklyException, WeeklyData
        validators.py       # Pre-run sanity checks (coordinates, missing fields)
    scheduling/
        eligibility.py      # Build e_{j,t,k,d} and e^H_{j,t,d} matrices
        milp.py             # Pyomo Layer-1 MILP (solve_scheduling)
        greedy.py           # Greedy fallback scheduler
    routing/
        clustering.py       # Geographic cluster construction (Algorithm 1)
        distance.py         # OSMnx / Haversine distance matrix builder
        milp.py             # Pyomo Layer-2 routing MILP
        sa.py               # Simulated Annealing solver
        greedy.py           # Constructive greedy heuristic
        bfs.py              # Brute-force / branch-and-bound
    postprocess/
        standby.py          # Standby job ranking
        readiness.py        # Area readiness computation
        output.py           # Route maps, text summaries, reason codes
    config.py               # Loads JSON/YAML config; merges Airtable overrides
    pipeline.py             # Top-level orchestrator (the 9-step pipeline)
\end{verbatim}

\paragraph{Design principles.}
\begin{enumerate}
    \item \textbf{Stateless stages.} Each stage is a pure function: it receives typed dataclass inputs and returns typed outputs. No global mutable state.
    \item \textbf{Fail-fast validation.} \texttt{validators.py} runs before the pipeline and rejects runs with missing columns, out-of-bounds coordinates, or conflicting exceptions. Errors are collected, not thrown one at a time.
    \item \textbf{Solver-agnostic interface.} Both \texttt{scheduling/milp.py} and \texttt{routing/milp.py} expose the same \texttt{solve(data, config) -> Result} signature. If Pyomo/CBC is unavailable, \texttt{pipeline.py} falls back to the greedy variant without code changes.
    \item \textbf{Idempotent runs.} Given the same inputs and config, the engine produces identical outputs. Random seeds for SA are drawn from the config file.
    \item \textbf{Separation of I/O.} File reads, API calls, and map rendering live in \texttt{data/loaders.py} and \texttt{postprocess/output.py}. Core logic never touches the filesystem directly.
\end{enumerate}

\paragraph{Pipeline orchestration (\texttt{pipeline.py}).}
The top-level function \texttt{run\_weekly\_plan(config\_path, data\_dir)} executes the nine stages listed in Section~\ref{sec:pipeline} in sequence.
Each stage logs its runtime and emits a structured summary (JSON) so that failures can be diagnosed without re-running the full pipeline.

\paragraph{Testing strategy.}
Unit tests cover individual stages with small synthetic fixtures.
An integration test runs the full pipeline on the example CSVs and asserts that:
\begin{itemize}
    \item every assigned job satisfies eligibility;
    \item no technician exceeds $T_{\max}$;
    \item helper jobs have exactly one primary and one distinct helper;
    \item output files are generated without errors.
\end{itemize}

% ===================================================================
\section*{Code vs.\ Configuration vs.\ Airtable Split}
\label{sec:split}
% ===================================================================

Not all system behaviour belongs in code.
The following table defines which concerns live in Python source, which in a version-controlled config file (JSON/YAML), and which in Airtable (editable by the owner at any time).

\medskip
{\small
\begin{longtable}{p{5.2cm}ccc}
\toprule
\textbf{Concern} & \textbf{Code} & \textbf{Config} & \textbf{Airtable} \\
\midrule
MILP formulation \& solver logic              & $\bullet$ &           &           \\
Clustering algorithm                          & $\bullet$ &           &           \\
SA / Greedy / BFS heuristics                  & $\bullet$ &           &           \\
Distance-matrix computation                   & $\bullet$ &           &           \\
Output rendering (maps, text)                 & $\bullet$ &           &           \\
\midrule
$\lambda$ (uniformity weight)                 &           & $\bullet$ &           \\
$T_{\max}$ (daily time budget)                &           & $\bullet$ &           \\
$P$ (global max stops)                        &           & $\bullet$ &           \\
$R$ (cluster radius)                          &           & $\bullet$ &           \\
$r$ (inter-stop radius)                       &           & $\bullet$ &           \\
Seasonal weight profiles $(w_g,w_f,w_a,w_r)$  &           & $\bullet$ &           \\
Solver time-out (seconds)                     &           & $\bullet$ &           \\
SA hyper-parameters ($T_{\text{start}}, \alpha, K$) &     & $\bullet$ &           \\
Random seed                                   &           & $\bullet$ &           \\
\midrule
Vehicle roster \& capability tags             &           &           & $\bullet$ \\
Vehicle availability days                     &           &           & $\bullet$ \\
Vehicle speed / cost / capacity overrides     &           &           & $\bullet$ \\
Technician roster \& skills                   &           &           & $\bullet$ \\
Technician home-depot coordinates             &           &           & $\bullet$ \\
Weekly exceptions (PTO, maintenance)          &           &           & $\bullet$ \\
Prohibited technician--job pairings           &           &           & $\bullet$ \\
Geographic area definitions ($\mathcal{A}$)   &           &           & $\bullet$ \\
Special-route flags (radiator, NH, helper)    &           &           & $\bullet$ \\
Job candidate pool (from ServiceM8 sync)      &           &           & $\bullet$ \\
\bottomrule
\end{longtable}
}

\paragraph{Rationale.}
\begin{itemize}
    \item \textbf{Code} holds logic that requires developer review to change---model structure, algorithm implementations, output formats. Changes go through version control and testing.
    \item \textbf{Config (JSON/YAML)} holds numerical parameters that a developer or technically confident owner adjusts between runs. The config file is version-controlled alongside the code so that every run is reproducible. A schema validator rejects unknown keys or out-of-range values before the pipeline starts.
    \item \textbf{Airtable} holds operational data that changes weekly or on demand---rosters, exceptions, pairings, area maps. The owner edits these freely; the engine pulls the latest state at the start of each run via the Airtable API. No code deployment is needed.
\end{itemize}

\paragraph{Override precedence.}
When the same parameter appears in both the config file and an Airtable field (e.g.\ a per-vehicle capacity override), the Airtable value takes precedence.
This lets the owner make one-off adjustments without touching the config file, while the config file provides sensible defaults.

\paragraph{Versioning and audit trail.}
Every planning run records:
\begin{enumerate}
    \item The config-file hash (git commit or content hash).
    \item A snapshot of the Airtable tables consumed (written to a timestamped JSON file).
    \item The engine version (Python package version or git SHA).
\end{enumerate}
This triple ensures that any past run can be fully reproduced or audited.

% ===================================================================
\section*{Validation and Testing Approach}
\label{sec:validation}
% ===================================================================

Correctness is enforced at three levels: input validation (before the pipeline runs), constraint verification (after each layer), and regression testing (during development).

\subsection*{Input Validation (Pre-run)}

\texttt{validators.py} performs the following checks when data is loaded. If any check fails, the run is aborted with a consolidated error report---not a single exception.

\begin{enumerate}
    \item \textbf{Schema completeness.} Every required CSV column is present and non-null for every row.
    \item \textbf{Coordinate bounds.} All latitudes fall within $[41, 48]$ and longitudes within $[-74, -67]$ (approximate New England bounding box). Out-of-range values indicate geocoding errors.
    \item \textbf{Referential integrity.} Every \texttt{scope\_id} in the exceptions file matches an existing technician or vehicle ID. Every \texttt{area\_id} on a job exists in the area definitions.
    \item \textbf{Availability consistency.} After applying exceptions, at least one (technician, vehicle, day) triple must remain feasible; otherwise the run would produce an empty schedule.
    \item \textbf{Duplicate detection.} Duplicate \texttt{job\_id}, \texttt{tech\_id}, or \texttt{vehicle\_id} values are flagged.
    \item \textbf{Config schema.} The JSON/YAML config is validated against a schema: all numerical parameters are present, positive, and within documented bounds (e.g.\ $0 \leq \lambda \leq 10\,000$, $R > 0$).
\end{enumerate}

\subsection*{Post-Layer Constraint Verification}

After Layer~1 and Layer~2 each complete, the engine runs an assertion suite on the solution before proceeding.

\paragraph{Layer~1 assertions.}
\begin{itemize}
    \item Every assigned job appears at most once across all (technician, vehicle, day) triples: $\sum_{t,k,d} y_{j,t,k,d} \leq 1$.
    \item Eligibility is respected: $y_{j,t,k,d} = 1 \Rightarrow e_{j,t,k,d} = 1$.
    \item Vehicle capacity is satisfied per slot: $\sum_j y_{j,t,k,d} \leq Q_k$.
    \item Service-time budget holds: $\sum_j \tau_j y_{j,t,k,d} \leq T_{\max}$.
    \item Each technician uses at most one vehicle per day: $\sum_k z_{t,k,d} \leq 1$.
    \item Each vehicle is used by at most one technician per day: $\sum_t z_{t,k,d} \leq 1$.
    \item Every helper job ($j \in \mathcal{J}^H$) that is assigned has exactly one helper: $\sum_t y^H_{j,t,d} = 1$ when $\sum_{t,k} y_{j,t,k,d} = 1$.
    \item The helper technician differs from the primary: $y^H_{j,t,d} + \sum_k y_{j,t,k,d} \leq 1$ for all $t$.
\end{itemize}

\paragraph{Layer~2 assertions (per route).}
\begin{itemize}
    \item Route departs from and returns to the depot (node~0).
    \item No node is visited more than once.
    \item Flow conservation holds at every node: in-degree equals out-degree.
    \item Total travel-plus-service time $\leq T_{\max}$.
    \item Number of visited stops $\leq Q_k$.
    \item No arc violates the inter-stop limit $r$.
    \item All visited nodes lie within cluster radius $R$ of the depot.
\end{itemize}

If any assertion fails, the pipeline logs the violation, marks the affected route as \texttt{INFEASIBLE}, and continues with the remaining bundles rather than aborting the entire run.

\subsection*{Regression and Unit Testing}

\paragraph{Unit tests.}
Each module exposes pure functions that can be tested in isolation:
\begin{itemize}
    \item \texttt{test\_eligibility}: synthetic 3-job, 2-tech, 2-vehicle scenario; verify $e_{j,t,k,d}$ matrix matches hand-computed values.
    \item \texttt{test\_scheduling\_milp}: small instance where the optimal assignment is known; compare solver output to expected.
    \item \texttt{test\_clustering}: verify cluster membership given a fixed radius and technician location.
    \item \texttt{test\_routing\_milp}: 4-node instance with known optimal tour; compare objective and route.
    \item \texttt{test\_sa}: seeded SA run on the same 4-node instance; verify feasibility and non-regression of objective.
    \item \texttt{test\_helper}: 2-tech, 1 helper job; confirm that exactly one helper is assigned and differs from the primary.
\end{itemize}

\paragraph{Integration test.}
A single end-to-end test runs the full pipeline on the example CSVs (\texttt{data/example\_*.csv}) and asserts:
\begin{itemize}
    \item no validation errors;
    \item at least one route is generated;
    \item all post-layer assertions pass;
    \item output files (maps, summaries, reason codes) are created and non-empty.
\end{itemize}

\paragraph{Golden-file comparison.}
For deterministic runs (MILP-only, or SA with a fixed seed), the integration test compares output against a checked-in golden file.
A mismatch triggers a diff report so that unintended behavioural changes are caught immediately.

% ===================================================================
\section*{Review Flags and Exception Outputs}
\label{sec:flags}
% ===================================================================

The engine produces a structured set of review flags alongside every weekly plan.
These flags are designed for the non-technical owner to scan quickly and decide whether to approve, adjust, or re-run.

\subsection*{Flag Categories}

Each flag has a severity (\texttt{INFO}, \texttt{WARN}, \texttt{CRITICAL}), a short code, and a human-readable message.

\medskip
{\small
\begin{longtable}{llp{7.5cm}}
\toprule
\textbf{Code} & \textbf{Severity} & \textbf{Meaning} \\
\midrule
\texttt{DUP\_ADDR}    & INFO     & Two or more jobs share the same address. Grouped for route evaluation; review for possible duplicates. \\
\texttt{CROSS\_AREA}  & INFO     & A route mixes jobs from two or more areas. May be optimal geographically but warrants a glance. \\
\texttt{WEAK\_STANDBY} & WARN    & Fewer than 2 feasible standby jobs exist for a route. If a scheduled job drops, coverage is thin. \\
\texttt{VEH\_BOTTLENECK} & WARN  & A vehicle is assigned at 100\,\% capacity on 4+ days. Little room for ad-hoc or emergency jobs. \\
\texttt{TECH\_OVERLOAD} & WARN   & A technician's total weekly service time exceeds 90\,\% of their available hours. \\
\texttt{HELPER\_TRAVEL} & WARN   & The helper technician's home depot is more than $R$ metres from the primary technician's depot. Travel time may be underestimated. \\
\texttt{ROUTE\_DROP}  & WARN     & One or more Layer-1-scheduled jobs were dropped during Layer-2 routing because travel time pushed the route past $T_{\max}$. \\
\texttt{GEOCODE\_OOB} & CRITICAL & A job's coordinates fall outside the New England bounding box. Likely a geocoding error; the job is excluded and must be corrected. \\
\texttt{NO\_FEASIBLE} & CRITICAL & A (technician, vehicle, day) bundle produced by Layer~1 has no feasible route at all in Layer~2. The entire bundle is unserved. \\
\texttt{SOLVER\_TIMEOUT} & WARN  & A solver (MILP or SA) hit its time limit before proving optimality. The returned solution is feasible but may not be optimal. \\
\bottomrule
\end{longtable}
}

\subsection*{Exception / Exclusion Report}

For every job in the weekly candidate pool that is \emph{not} assigned to a route, the engine emits a reason code explaining why.
This report is the primary traceability artefact.

\medskip
{\small
\begin{longtable}{lp{9cm}}
\toprule
\textbf{Reason Code} & \textbf{Explanation} \\
\midrule
\texttt{ALREADY\_ASSIGNED}   & Job was assigned in a prior run and has not been reset. \\
\texttt{STATUS\_EXCLUDED}    & Job status is \texttt{excluded} or \texttt{cancelled} in ServiceM8. \\
\texttt{NO\_ELIGIBLE\_TRIPLE}& No (technician, vehicle, day) triple satisfies all eligibility conditions ($e_{j,t,k,d} = 0$ everywhere). \\
\texttt{CAPACITY\_FULL}      & All eligible slots were already at vehicle capacity $Q_k$. \\
\texttt{TIME\_BUDGET}        & Including this job would exceed the service-time budget $T_{\max}$ in every eligible slot. \\
\texttt{CLUSTER\_RADIUS}     & Job is outside the cluster radius $R$ from every available technician's depot. \\
\texttt{SKILL\_MISMATCH}     & No available technician has the required skills. \\
\texttt{VEHICLE\_MISMATCH}   & No available vehicle has the required capability tags (e.g.\ radiator job but V-2 unavailable). \\
\texttt{HELPER\_UNAVAIL}     & Job requires a helper ($h_j = 1$) but no second technician is eligible on any feasible day. \\
\texttt{PROHIBITED\_PAIR}    & All eligible technicians are on the prohibited-pairing list for this job. \\
\texttt{ROUTE\_DROP}         & Job was scheduled by Layer~1 but dropped during Layer~2 routing due to travel-time infeasibility. \\
\bottomrule
\end{longtable}
}

\subsection*{Output Format}

Flags and exclusion reasons are written to two files per run:

\begin{itemize}
    \item \texttt{flags\_\{date\}.json} --- array of flag objects, each with \texttt{code}, \texttt{severity}, \texttt{message}, and references to the affected route, technician, or job.
    \item \texttt{exclusions\_\{date\}.json} --- array of exclusion objects, each with \texttt{job\_id}, \texttt{reason\_code}, and a human-readable \texttt{detail} string.
\end{itemize}

Both files are included in the review packet sent to the owner alongside route maps and summaries.
\texttt{CRITICAL} flags are additionally surfaced as bold-red banners at the top of the summary output, ensuring they are not overlooked.

% ===================================================================
\section*{Future Considerations}
% ===================================================================

\begin{itemize}
    \item Iterative feedback loop between Layer~1 and Layer~2 (re-schedule infeasible bundles).
    \item Exact stop sequencing and appointment--time optimization (Version~2).
    \item Integration with real-time traffic and drive-time APIs.
    \item Historical job-success and no-show rates to refine scoring weights.
    \item Multi-week rolling horizon planning.
    \item Machine-learning--assisted demand forecasting for proactive area readiness.
\end{itemize}

\section*{ SW solution considerations}

The reading of the input (data.txt):  
\begin{itemize}
\item 
Parses the data file (handling the comma-decimal format in numbers).
\item Converts latitude and longitude to float.
\item  Uses the Google Maps Distance Matrix API or OSMnx's package to compute the matrix \(d_{i,j}\) of real-world driving distances between all pairs of locations.
\end{itemize}

Get the starting location and calculate the checkpoints inside the corresponding clusters.

For each of the transport type (for now it is Foot and Car) calculate distances (meters) as time(minutes) between checkpoints and between starting locations and checkpoints.

Approaches to find the best routes:
\begin{itemize}
    \item Use dynamic programming approach (Bellman-Ford type) to find the max revenue trip
    \item Use brute-force with branch-and-bound for every possible set (using some heuristics)
    \item Solve MILP Optimization problem mentioned
    \item Use Simulated Annealing approach
    \item Heuristics
\end{itemize}

The resulting solution route is displayed on the map.

\section*{Finding the best route}

\subsection*{MILP Optimization problem}

We can directly solve the described optimization problem using packages like Pyomo with open solvers (for example, CBC or GLPK).

The overall worst-case number of variables is \(N*N + N = N(N+1)\), where \(N\) is within \(100-3300\), so the order of the number of variables is \(100000 =10^5\), looks like it is too large even for commercial solvers.

However, we may note that we don't have to use the whole set of checkpoints, but only those that are situated within the maximum zone distance radius. For file data.txt, we have that each such zone gives only at max 100 checkpoints. Therefore, we have an order of $10^3$ of variables that is realistic to solve with open solvers.

\subsection*{Heuristics}

The possible approach for heuristic is the following: from the starting location we list the points within accessing radius, and sort them by the value of revenue of the point minus the cost of the way to it. We select the maximal value for this sorting and choose is as the first nexus (route). 

After this, we start from this selected point and select new route with the same sorting, always considering that new checkpoint must be inside accessing radius for the start and the current point.

We continue until we either have \(P\) checkpoints selected or we cannot find the point where time for the whole route is less than \(T_{max}\).


\subsection*{Brute-force branch-and-bound}

We implement the BFS approach: 

For the starting point we calculate for all accessible checkpoints the price of it as the value of revenue of the point minus the cost of the way to it, as well as the time to reach this points. These points are put in the queue. Then, as in BFS we took those points from the queue and calculate again the corresponding neighbors (that are not already in the routes) and store them in the queue. 

We proceed with this procedure until we either get P length route or the total time constraint cannot be satisfied.  

Also, we use branch and bound approach. As far as we get heuristics solution and starting minima value we reject options where the revenue be certainly less than already calculated.

The overall complexity/memory of this approach is \(O(N^P)\),where \(N\) is the number of nodes in accessible zone, i.e. we have \(100^{14}=10^{28}\).

\subsection*{Simulated annealing}


Simulated Annealing (SA) is a metaheuristic inspired by the annealing process in metallurgy. It is suitable for combinatorial optimization problems such as our inspection routing task, where the objective is to select and order at most $P$ nodes to maximize revenue while satisfying time, distance, and capacity constraints.

\subsubsection*{Key Concepts}

\begin{itemize}
    \item \textbf{State:} A feasible inspection route starting and ending at node $0$, visiting a subset of nodes $\Omega \subseteq \{1, \ldots, N\}$.
    \item \textbf{Objective Function:}
    \[
    \max\left( \sum_{i=0}^{N}\sum_{j=0}^{N} w_{i} x_{i,j} - \sum_{i=0}^{N}\sum_{j=0}^{N} c_k \cdot d_{i,j} \cdot x_{i,j} \right)
    \]
    \item \textbf{Constraints:} Each candidate route must satisfy the vehicle capacity limit $|\Omega| \leq Q_k$, total time limit $T_{\text{max}}$, and local and global distance constraints.
    \item \textbf{Temperature:} Controls the acceptance of worse solutions to escape local minima.
    \item \textbf{Cooling schedule:} Determines how temperature decreases over time.
\end{itemize}

\subsubsection*{Neighborhood Operations}

To explore the solution space, we define several perturbation operators:

\begin{itemize}
    \item \textbf{Node Insertion:} Randomly insert a new node $j \notin \Omega$ into the current route at a random position, if it satisfies constraints.
    \item \textbf{Node Deletion:} Randomly remove a node $j \in \Omega$, except the start node.
    \item \textbf{Node Swap:} Swap two nodes $i, j \in \Omega$, modifying the order of visits.
    \item \textbf{External Swap:} Replace a node $i \in \Omega$ with another node $j \notin \Omega$.
    \item \textbf{2-opt Flip:} Reverse a subsegment of the route to reduce total distance.
\end{itemize}

\subsubsection*{Simulated Annealing Algorithm}

\begin{algorithm}[H]
\caption{Simulated Annealing for Inspection Routing}
\begin{algorithmic}[1]
\State Initialize route $\Omega_0$ with feasible nodes
\State $T \gets T_{\text{start}},\ T_{\text{end}} \gets$ final temperature
\While{$T > T_{\text{end}}$}
    \For{$k = 1$ to $K$ (inner loop)}
        \State $r_{\text{new}} \gets$ apply a random neighbor operation to $\Omega$
        \If{$r_{\text{new}}$ is feasible}
            \State $\Delta \gets \text{Objective}(r_{\text{new}}) - \text{Objective}(\Omega)$
            \If{$\Delta > 0$}
                \State Accept new route
            \ElsIf{$\exp\left(\frac{\Delta}{T}\right) > \text{random}(0,1)$}
                \State Accept new route with probability
            \EndIf
        \EndIf
    \EndFor
    \State $T \gets \alpha T$ \Comment{Cooling: $\alpha \in (0,1)$}
\EndWhile
\end{algorithmic}
\end{algorithm}

\subsubsection*{Feasibility Check}

Every new route $\Omega'$ generated during neighborhood exploration is accepted only if:
\begin{itemize}
    \item The number of nodes visited does not exceed $Q_k$.
    \item The total travel and service time is within $T_{\text{max}}$:
    \[
    \sum_{(i,j) \in \Omega'} \frac{d_{i,j}}{v_k} + \tau_j \leq T_{\text{max}}
    \]
    \item The distances $d_{0,j} \leq R$ for entering and exiting the cluster.
    \item The step distances $d_{i,j} \leq r$ between consecutive visits.
\end{itemize}

\subsubsection*{Advantages}

Simulated Annealing avoids local optima by probabilistically accepting worse solutions during early stages. With well-chosen operators and cooling schedules, SA efficiently explores complex route spaces.

\subsubsection*{Implementation Notes}

\begin{itemize}
    \item The initial solution can be constructed greedily by selecting the $P$ nodes with highest $(w_i - c \cdot d_{0i})$ ratios.
    \item Use a priority queue or randomized selection for neighbor generation.
    \item Recompute route cost incrementally where possible to improve efficiency.
\end{itemize}







\end{document}
